\documentclass[11pt]{article}

\usepackage[margin=1.08in]{geometry}
\usepackage{amsmath,amssymb,amsthm}
\usepackage{enumitem}
\usepackage[hidelinks]{hyperref}

\newtheorem{theorem}{Theorem}[section]
\newtheorem{proposition}[theorem]{Proposition}
\newtheorem{lemma}[theorem]{Lemma}
\newtheorem{corollary}[theorem]{Corollary}
\theoremstyle{definition}
\newtheorem{definition}[theorem]{Definition}
\theoremstyle{remark}
\newtheorem{remark}[theorem]{Remark}

\newcommand{\C}{\mathbb C}
\newcommand{\A}{\mathbb A}
\newcommand{\KX}{\mathcal K_X}
\newcommand{\KS}{\mathcal K_S}
\newcommand{\ord}{\operatorname{ord}}
\newcommand{\wt}{\operatorname{wt}}
\newcommand{\Jac}{\operatorname{Jac}}
\newcommand{\Res}{\operatorname{Res}}

\title{Squarefree Common-Root Boundaries\\
for Reciprocal Keller Pairs}
\author{Atharva Vaidya}
\date{26 July 2026}

\begin{document}
\maketitle

\begin{abstract}
Let \(g\geq2\), let \(1<a<b\) be coprime, and put
\(n=ga\), \(m=gb\), and \(N=n+m-2\).
We consider reciprocal polynomial pairs \(P,Q\in\C[X,\tau]\) satisfying
\[
 \deg_X[\tau^j]P\leq n-j,\qquad
 \deg_X[\tau^j]Q\leq m-j
\]
and the homogenized Keller identity
\[
 \tau(P_XQ_\tau-P_\tau Q_X)+nPQ_X-mQP_X
 =-c\tau^N,\qquad c\in\C^\times.
\]
We prove that such a pair cannot have boundary values
\[
 P(X,0)=R(X)^a,\qquad Q(X,0)=R(X)^b
\]
when \(R\) is a monic squarefree polynomial of degree \(g\).

The proof first shows that every strict first compact Newton face
encountered by the least-residual induction at a root of \(R\) is
automatically matched in the two coordinates and is a reduced common-root
shift.  These shifts glue through all orders below \(g\).
A least-residual induction and the reciprocal degree bounds then force
the exact resonant normal form
\[
 P=S^aA(\tau^g/S),\qquad Q=S^bB(\tau^g/S),
\]
where \(S\) is a reciprocal polynomial deformation of \(R\) and \(A,B\)
are arbitrary polynomials of the allowed degrees.  The homogenized
Keller bracket of every such pair is identically zero, giving the
contradiction.

This is an exclusion theorem for one boundary architecture.  It does
not prove the plane Jacobian conjecture.  In particular, repeated-root
boundaries admit local factor-splitting and Kummer modes that are not
covered by the argument.
\end{abstract}

\section{Introduction and scope}

The two-dimensional Jacobian conjecture asks whether every polynomial
map \(f=(p,q):\A^2_\C\to\A^2_\C\) with
\(\Jac(p,q)\in\C^\times\) is a polynomial automorphism.  The conjecture
remains open.  A classical approach studies the configuration of roots
and Newton data at infinity; see Biggers--Moh--Fried
\cite{BiggersMohFried} and the subsequent standard-pair and admissible
chain developments of Guccione, Guccione, Horruitiner, and Valqui
\cite{GGVSystem,GGVLowerSide,GGHVAlgorithms,GGHV2022}.

This paper isolates a boundary condition that arises naturally in a
reciprocal chart at infinity.  Suppose the top boundary polynomials are
coprime powers of one degree-\(g\) polynomial:
\[
 P(X,0)=R(X)^a,\qquad Q(X,0)=R(X)^b,
 \qquad \gcd(a,b)=1.
\]
Our result says that \(R\) cannot be squarefree when the entire
reciprocal polynomial pair satisfies the exact Keller identity.
Equivalently, every candidate that reaches this common-power boundary
must retain a repeated root.

There are three reasons that the conclusion is not merely a local
Newton-polygon observation.  First, the matching \(Q\)-face has to be
deduced from the actual polynomial Keller identity; it cannot safely be
read term by term from a Laurent standard presentation.  Second, local
root shifts must glue into one polynomial approximate root while
respecting the reciprocal degree bounds.  Third, the slope-\(g\)
resonance does not force \(P\) and \(Q\) to remain relative powers.
The proof retains independent resonant polynomials \(A\) and \(B\) and
shows that the whole resulting cone nevertheless has zero bracket.

\paragraph{Relation to earlier frameworks.}
Guccione--Guccione--Valqui proved that failure of the Jacobian
conjecture is equivalent to solvability of a finite system of polynomial
equations \cite{GGVSystem}; later work organizes possible standard
minimal pairs and Newton corners algorithmically
\cite{GGVLowerSide,GGHVAlgorithms,GGHV2022}.  The theorem below is
formulated directly for the reciprocal polynomial Keller identity.
It does not reprove the reduction from an arbitrary counterexample to
any particular standard system, and it does not assert that every
common-root polynomial arising in every normalization is squarefree.

Makar-Limanov and Trakhtenberg study Newton-resolution integrality and
polynomiality conditions for a Jacobian mate, including a decreasing
nonprincipal power index \cite{MakarLimanovTrakhtenberg}.  Their
resolution constraints are close in spirit to the approximate-root
analysis here, but we do not import their results, and we are not aware
of the squarefree reciprocal-boundary exclusion in their work.
Shaska proves that every \(\C^\times\)-equivariant Keller self-map of
\(\A^2\) is an automorphism \cite{Shaska}.  Our final resonant cone has
an exact weighted form, but the present theorem is neither an
equivariant-map theorem nor a consequence claimed from that result.
These comparisons are statements of scope, not claims of priority.

\paragraph{Logical boundary.}
The theorem does not prove or disprove the plane Jacobian conjecture.
It excludes the squarefree case of a specific reciprocal common-root
boundary.  Repeated roots are materially different: the local
coordinate is no longer \(R\) itself, a strict face may split a repeated
factor, and sub-\(g\) Kummer kernel modes can occur.  Section
\ref{sec:repeated} records precisely where the present proof stops.

\section{Reciprocal Keller pairs}

\subsection{The reciprocal chart}

Let \(p,q\in\C[x,y]\) have degree at most \(n,m\), respectively.  Put
\[
 X=\frac{x}{y},\qquad \tau=\frac1y
\]
and define
\begin{equation}
 P(X,\tau)=\tau^n p(X/\tau,1/\tau),\qquad
 Q(X,\tau)=\tau^m q(X/\tau,1/\tau).
\label{eq:reciprocal}
\end{equation}
These are polynomials.  If
\[
 P=\sum_{j\geq0}\tau^jP_j(X),\qquad
 Q=\sum_{j\geq0}\tau^jQ_j(X),
\]
then
\begin{equation}
 \deg P_j\leq n-j,\qquad
 \deg Q_j\leq m-j.
\label{eq:degree-bounds}
\end{equation}
We use the convention \(\deg0=-\infty\), so the bounds also imply
\(P_j=0\) for \(j>n\) and \(Q_j=0\) for \(j>m\).

\begin{definition}
For positive integers \(n,m\), define the homogenized Keller bracket
\begin{equation}
 \KX(P,Q)
 :=
 \tau(P_XQ_\tau-P_\tau Q_X)+nPQ_X-mQP_X.
\label{eq:kx}
\end{equation}
A \emph{reciprocal Keller pair of bidegree \((n,m)\)} is a pair
\(P,Q\in\C[X,\tau]\) satisfying \eqref{eq:degree-bounds} and
\[
 \KX(P,Q)=-c\tau^{n+m-2}
\]
for some \(c\in\C^\times\).
\end{definition}

\begin{proposition}[Reciprocal covariance of the Jacobian]
\label{prop:reciprocal}
For \eqref{eq:reciprocal},
\[
 \Jac_{x,y}(p,q)
 =-\tau^{-n-m+2}\KX(P,Q).
\]
In particular, \(\Jac(p,q)=c\) if and only if
\[
 \KX(P,Q)=-c\tau^{n+m-2}.
\]
Conversely, the degree bounds make
\[
 p(x,y)=y^nP(x/y,1/y),\qquad
 q(x,y)=y^mQ(x/y,1/y)
\]
polynomials, so every reciprocal Keller pair dehomogenizes to a plane
Keller pair.
\end{proposition}

\begin{proof}
The coordinate Jacobian is
\[
 \frac{\partial(X,\tau)}{\partial(x,y)}=-\tau^3.
\]
Moreover,
\[
 \frac{\partial(p,q)}{\partial(X,\tau)}
 =\tau^{-n-m-1}\KX(P,Q).
\]
Multiplying these determinants proves the formula.  The converse
polynomiality is exactly \eqref{eq:degree-bounds}.
\end{proof}

\subsection{Statement of the theorem}

\begin{theorem}[Squarefree common-root exclusion]
\label{thm:main}
Let
\[
 g\geq2,\qquad 1<a<b,\qquad \gcd(a,b)=1,
\]
and put
\[
 n=ga,\qquad m=gb,\qquad N=n+m-2.
\]
There is no reciprocal Keller pair \(P,Q\in\C[X,\tau]\) satisfying
\begin{equation}
 P(X,0)=R(X)^a,\qquad Q(X,0)=R(X)^b
\label{eq:boundary}
\end{equation}
for a monic squarefree polynomial \(R\in\C[X]\) of degree \(g\).
\end{theorem}

The proof gives a more precise normal-form statement.

\begin{theorem}[Forced resonant normal form]
\label{thm:normal-form}
Under the hypotheses of Theorem~\ref{thm:main}, the Newton-face and
reciprocal-degree constraints force polynomials
\[
 S(X,\tau)
 =R(X)+\sum_{1\leq e<g}\tau^eL_e(X),
 \qquad \deg L_e\leq g-e,
\]
and
\[
 A(z)=\sum_{q=0}^{a}A_qz^q,\qquad
 B(z)=\sum_{q=0}^{b}B_qz^q,\qquad A_0=B_0=1,
\]
such that
\begin{equation}
 \boxed{
 P=S^aA(\tau^g/S),\qquad Q=S^bB(\tau^g/S).
 }
\label{eq:normal-form}
\end{equation}
For arbitrary \(A,B\) in this formula,
\[
 \KX(P,Q)=0.
\]
\end{theorem}

Thus Theorem~\ref{thm:normal-form} is the structural statement obtained
immediately before the contradiction.  At each face used in its proof,
vanishing below the forcing weight follows from the exact identity
\(\KX(P,Q)=-c\tau^N\).

\section{Compact Newton faces}
\label{sec:faces}

\subsection{A binomial warm-up}

The following computation motivates both the physical slope \(g\) and
the need to retain every monomial on a compact face.

\begin{lemma}[An isolated binomial face cannot supply the scalar]
\label{lem:binomial}
Write
\[
 b=aq+r,\qquad 1\leq r<a,
\]
and consider the local bracket
\[
 \KS(F,G)
 =\tau(F_sG_\tau-F_\tau G_s)+gaFG_s-gbGF_s.
\]
If the isolated \(P\)-face is
\[
 F=s^a+c_0\tau^d,\qquad c_0\neq0,\qquad d<ga,
\]
then cancellation forces the polynomial \(Q\)-chain
\[
 G_0=\sum_{j=0}^{q}C_j\tau^{dj}s^{b-aj},
 \qquad
 C_j=c_0^j\binom{b/a}{j}.
\]
Its uncancelled endpoint is
\begin{equation}
 c_0r(ga-d)C_q\,\tau^{d(q+1)}s^{r-1}.
\label{eq:binomial-endpoint}
\end{equation}
No distinct polynomial \(Q\)-monomial cancels this endpoint.
It cannot be the nonzero scalar forcing at order
\(N=g(a+b)-2\).
\end{lemma}

\begin{proof}
For \(G_j=C_j\tau^{dj}s^{b-aj}\), the \(s^a\)-part and
\(c_0\tau^d\)-part of the bracket are
\[
 aj(d-ga)C_j\tau^{dj}s^{b-a(j-1)-1}
\]
and
\[
 c_0(ga-d)(b-aj)C_j
 \tau^{d(j+1)}s^{b-aj-1},
\]
respectively.  Consecutive supports coincide, forcing
\[
 C_{j+1}
 =c_0\frac{b-aj}{a(j+1)}C_j.
\]
This gives the displayed binomial coefficients and leaves
\eqref{eq:binomial-endpoint}.  A further monomial
\(h\tau^es^k\) could meet the endpoint only with
\[
 (e,k)=(d(q+1),r-a)
 \quad\hbox{or}\quad
 (e,k)=(dq,r).
\]
The first has \(k<0\), while the second is the already fixed last
chain monomial.

For the endpoint to be scalar in \(s\), one needs \(r=1\).  Then
\[
 N=g(a+b)-2=ga(q+1)+g-2.
\]
The equality \(d(q+1)=N\), together with \(d\leq ga\) and \(g\geq2\),
forces \(g=2\) and \(d=ga\).  But then the coefficient
\(ga-d\) is zero.  At \(d=ga\) the recurrence itself is not forced:
both displayed bracket contributions vanish term by term, and the face
still supplies no nonzero scalar.
\end{proof}

\begin{remark}
Lemma~\ref{lem:binomial} assumes that there are no interior
\(P\)-monomials on the edge.  The main argument cannot make that
assumption.  We now treat the entire compact face and prove that the
matching \(Q\)-face is automatic for an actual reciprocal Keller pair.
\end{remark}

\subsection{Local coordinates and covariance}

Let \(\alpha\) be a simple root of \(R\).  During the induction we use a
polynomial approximate root
\[
 s=S(X,\tau),\qquad S(X,0)=R(X).
\]
Since \(R'(\alpha)\neq0\), the formal implicit function theorem gives a
local inverse \(X=\chi(s,\tau)\).  This change preserves the least
\(\tau\)-order and the vanishing order of its first coefficient:
\[
 \tau^jT_j(X)+O(\tau^{j+1})
 =
 \tau^jT_j(\chi(s,0))+O(\tau^{j+1}).
\]
The inverse substitution places \(P\) and \(Q\) in
\(\C[[s,\tau]]\).  Because all Newton weights used below are positive,
each initial face is finite and has nonnegative exponents; hence the local
lattice-coset decompositions are legitimate in this formal series ring.

\begin{lemma}[Covariance under a moving root coordinate]
\label{lem:covariance}
If \(P=f(S,\tau)\) and \(Q=k(S,\tau)\), then
\begin{equation}
 \KX(P,Q)=S_X\KS(f,k),
\label{eq:covariance}
\end{equation}
where
\[
 \KS(f,k)
 =\tau(f_Sk_\tau-f_\tau k_S)+nfk_S-mkf_S
\]
and the \(\tau\)-derivatives on the right are taken with \(S\) fixed.
\end{lemma}

\begin{proof}
Use
\[
 P_X=S_Xf_S,\qquad
 P_\tau|_X=f_\tau|_S+S_\tau f_S
\]
and the analogous formulas for \(Q\).  The two terms containing
\(S_\tau f_Sk_S\) cancel.
\end{proof}

\subsection{The automatic matched-face theorem}

Give \(s,\tau\) primitive positive weights
\[
 \wt(s)=d',\qquad \wt(\tau)=h',\qquad
 \gcd(d',h')=1.
\]
The slope is called \emph{strict} when
\begin{equation}
 \frac{d'}{h'}<g.
\label{eq:strict}
\end{equation}
Put
\[
 z=\frac{\tau^{d'}}{s^{h'}}.
\]
Suppose the entire first compact \(P\)-face from \(s^a\) is
\begin{equation}
 P_0=s^aA(z),\qquad A(0)=1.
\label{eq:pface}
\end{equation}
Polynomiality means that every term \(z^j\) in \(A\) satisfies
\(h'j\leq a\).  By Lemma~\ref{lem:covariance}, the local Keller
equation has right side
\[
 -\frac{c}{S_X}\tau^N.
\]
The factor \(1/S_X\) is a formal local unit with nonzero constant term,
so the initial weight of the forcing remains \(h'N\).

\begin{lemma}[General \(Q\)-coset formula]
\label{lem:coset}
Let
\[
 Q_{u,v}=\tau^us^vB(z),\qquad B(0)\neq0,
\]
and put
\[
 w=h'u+d'v,\qquad
 \Delta=d'-gh',\qquad
 L=u+gv-gb.
\]
Then
\begin{align}
 \KS(P_0,Q_{u,v})
 =\tau^us^{a+v-1}\bigl(
 &aL\,AB+a\Delta z\,AB'\notag\\
 &+(gbh'-w)z\,A'B
 \bigr).
\label{eq:coset}
\end{align}
\end{lemma}

\begin{proof}
Differentiate \eqref{eq:pface} and \(Q_{u,v}\), using
\[
 s z_s=-h'z,\qquad \tau z_\tau=d'z,
\]
and collect the three Euler terms.  The coefficient independent of
\(z\) inside the parentheses is \(aL A(0)B(0)\).
\end{proof}

\begin{proposition}[Strict compact-face bridge]
\label{prop:strict-face}
For an actual reciprocal Keller pair satisfying \eqref{eq:boundary},
the lowest \(Q\)-face paired with the strict first \(P\)-face
\eqref{eq:pface} is automatically
\[
 Q_0=s^bB(z),\qquad B(0)=1.
\]
Moreover, every nontrivial such strict first compact face is a common
reduced shift:
\begin{equation}
 \boxed{
 P_0=(s+\kappa\tau^\delta)^a,\qquad
 Q_0=(s+\kappa\tau^\delta)^b,
 \quad
 \kappa\in\C^\times,\quad 1\leq\delta<g.
 }
\label{eq:reduced-face}
\end{equation}
The same conclusion holds if the first strict face is detected from
\(Q\).
\end{proposition}

\begin{proof}
Suppose the actual \(Q\) had a lowest face of weight \(w<bd'\).
Its bracket with \(P_0\) has weight
\[
 (a-1)d'+w<(a+b-1)d'.
\]
On the other hand, the right side of the Keller identity has weight
\(h'N\), and
\begin{align}
 h'N-(a+b-1)d'
 &=(a+b-1)(gh'-d')+h'(g-2)>0.
\label{eq:weight-gap}
\end{align}
The inequality remains strict when \(g=2\).  Thus the lower face must
vanish.  The constant term in \eqref{eq:coset} forces \(L=0\), so
\(u=g(b-v)\).  Since \(u,v\geq0\), one has \(v\leq b\), and hence
\[
 w=h'g(b-v)+d'v
 =bd'+(b-v)(gh'-d')\geq bd',
\]
a contradiction.

Therefore no \(Q\)-face lies below \(s^b\).  Every integral lattice
point on the weight-\(bd'\) line lies in the single coset through
\((u,v)=(0,b)\), because \(\gcd(d',h')=1\).  The complete face is
therefore \(Q_0=s^bB(z)\) with \(B(0)=1\).  Terms on adjacent faces
have strictly higher weight and do not enter the initial bracket.

Putting \(u=0,v=b,w=bd'\) in \eqref{eq:coset} yields
\begin{equation}
 \KS(P_0,Q_0)
 =(d'-gh')s^{a+b-1}z
 \bigl(aAB'-bA'B\bigr).
\label{eq:matched-bracket}
\end{equation}
Its weight is \((a+b-1)d'\), strictly below the forcing by
\eqref{eq:weight-gap}.  Lemma~\ref{lem:covariance} and the fact that
\(S_X\) is a local unit show that
\[
 aAB'-bA'B=0.
\]
If one of \(A,B\) were constant, the other would also be constant.
For a nontrivial face, differentiation of \(B^a/A^b\), normalization
at \(z=0\), and unique factorization in \(\C[z]\) give
\[
 B^a=A^b,\qquad A=C^a,\qquad B=C^b,\qquad C(0)=1.
\]

Let \(\ell=\deg A\).  Because \((d',h')\) is primitive, \(\ell\) is
the lattice length of the \(P\)-edge, and its horizontal drop is
\(h'\ell\leq a\).  From \(A=C^a\) one has \(a\mid\ell\).  Hence
\[
 a\leq\ell\leq h'\ell\leq a.
\]
Thus \(\ell=a\), \(h'=1\), and \(\deg C=1\).  Write
\(C=1+\kappa z\).  Since the face is strict, \(d'<g\).
Taking \(\delta=d'\) gives \eqref{eq:reduced-face}.
Interchanging \((P,n,a)\) with \((Q,m,b)\) changes the bracket only by
a sign, so the same argument applies when the strict face is first
detected from \(Q\).
\end{proof}

\begin{remark}[Why the full polynomial identity matters]
The conclusion of Proposition~\ref{prop:strict-face} is about the
lowest face of the actual polynomial \(Q\).  In a Laurent or Puiseux
standard presentation, polynomial projection can change finite-root
Taylor order.  Thus one cannot prove automatic matching by assigning a
weight to each unprojected standard term separately.  Formula
\eqref{eq:coset}, the scalar weight gap, and polynomiality are the
bridge.
\end{remark}

\section{Gluing through the subresonant orders}
\label{sec:subresonant}

We first absorb all common reduced shifts of order below \(g\).
Starting with \(S=R\), assume inductively that lower orders have been
removed and let \(0<j<g\) be the least order at which either residual
is nonzero:
\begin{equation}
\begin{aligned}
 P-S^a&=\tau^jT_j(X)+O(\tau^{j+1}),\\
 Q-S^b&=\tau^jU_j(X)+O(\tau^{j+1}).
\end{aligned}
\label{eq:subresidual}
\end{equation}

\begin{lemma}[A subresonant layer contains a strict root]
\label{lem:strict-root}
If \(T_j\neq0\), then at some root \(\alpha\) of \(R\), the point
defined by \(\tau^jT_j\) has slope strictly below \(g\).
The same holds for \(U_j\neq0\).
\end{lemma}

\begin{proof}
For each of the \(g\) distinct roots of \(R\), put
\[
 \sigma_\alpha
 =\min\{\ord_\alpha T_j,a\},\qquad
 h_\alpha=a-\sigma_\alpha.
\]
Squarefreeness gives
\[
 \sum_\alpha\sigma_\alpha\leq\deg T_j\leq ga-j,
\]
and hence
\begin{equation}
 j\leq\sum_\alpha h_\alpha.
\label{eq:degree-average}
\end{equation}
If no root were strict, then \(j\geq gh_\alpha\) at every root, so
\(\sum h_\alpha\leq j\).  Equality would hold everywhere and force
\(h_\alpha=j/g\) for all \(\alpha\), impossible because
\(0<j<g\).  Replace \(a\) by \(b\) for \(U_j\).
\end{proof}

\begin{proposition}[Subresonant common-root gluing]
\label{prop:subgluing}
At the order \(j\) in \eqref{eq:subresidual}, there is a polynomial
\(L_j\in\C[X]\) with
\[
 \deg L_j\leq g-j
\]
such that
\begin{equation}
 T_j=aR^{a-1}L_j,\qquad
 U_j=bR^{b-1}L_j.
\label{eq:glued-layer}
\end{equation}
Consequently the replacement
\[
 S\longmapsto S+\tau^jL_j
\]
removes both residuals at order \(j\) and preserves all reciprocal
degree bounds.
\end{proposition}

\begin{proof}
Suppose first that \(T_j\neq0\).  At a root where
\(\ord_\alpha T_j<a\), put \(h=a-\ord_\alpha T_j\).
Since \(h\geq1\) and \(j<g\), the supplied point has slope
\(j/h\leq j<g\).
The actual first face has no greater slope, so
Proposition~\ref{prop:strict-face} makes it a reduced shift with
primitive order \(\delta<g\).  Its expansion contains
\[
 a\kappa\tau^\delta s^{a-1}.
\]
Since \(j\) is the least residual order, \(\delta\geq j\).  Since the
face slope is \(\delta\),
\[
 \delta\leq\frac{j}{h}\leq j.
\]
Thus \(\delta=j\), \(h=1\), and
\[
 \ord_\alpha T_j=a-1
\]
at every active root.  At an inactive root the order is at least
\(a\).  Lemma~\ref{lem:strict-root} guarantees that this reasoning is
not vacuous.  It follows that
\[
 T_j=aR^{a-1}L_j.
\]
The degree bound gives
\[
 \deg L_j
 \leq(ga-j)-g(a-1)=g-j.
\]

At an active root, the common face identifies the leading
order-\(j\) coefficient of \(Q\) with
\(bL_j(\alpha)s^{b-1}\), modulo \(s^b\).  At an inactive root,
\(L_j(\alpha)=0\).  If \(\ord_\alpha U_j<b\) there, the \(Q\)-point
would expose a strict common face and make the root active in \(P\),
a contradiction.  Hence
\[
 U_j-bR^{b-1}L_j
\]
is divisible by \(R^b\).  Both terms have degree at most \(m-j<m\),
so the difference vanishes.

If \(T_j=0\) and \(U_j\neq0\), Lemma~\ref{lem:strict-root} produces a
strict \(Q\)-face.  Proposition~\ref{prop:strict-face} makes it common,
which would produce a nonzero \(P\)-coefficient at or before order
\(j\), contrary to \eqref{eq:subresidual}.  Thus a \(Q\)-only
subresonant layer is impossible.

Finally, the coefficient bound on \(L_j\) implies that every binomial
term in \((S+\tau^jL_j)^a\) and
\((S+\tau^jL_j)^b\) satisfies the corresponding reciprocal degree
bound.
\end{proof}

\begin{corollary}[Common-power form modulo \(\tau^g\)]
\label{cor:modg}
There is a polynomial
\begin{equation}
 S=R+\sum_{1\leq e<g}\tau^eL_e(X),
 \qquad \deg L_e\leq g-e,
\label{eq:S}
\end{equation}
such that
\begin{equation}
 P\equiv S^a\pmod{\tau^g},\qquad
 Q\equiv S^b\pmod{\tau^g}.
\label{eq:modg}
\end{equation}
\end{corollary}

\begin{proof}
Apply Proposition~\ref{prop:subgluing} successively to the least
remaining order below \(g\).  There are only \(g-1\) such orders.
\end{proof}

\section{The least-residual induction after resonance}
\label{sec:resonance}

The slope-\(g\) bracket prefactor in
\eqref{eq:matched-bracket} vanishes, so the resonant coefficients of
\(P\) and \(Q\) must remain independent.  Introduce the model
\begin{equation}
\begin{aligned}
 P_{\mathrm{mod}}
 &=\sum_{q=0}^{a}A_q\tau^{gq}S^{a-q},\\
 Q_{\mathrm{mod}}
 &=\sum_{q=0}^{b}B_q\tau^{gq}S^{b-q},
\end{aligned}
\qquad A_0=B_0=1.
\label{eq:model}
\end{equation}
Initially only \(A_0,B_0\) are chosen.  Corollary~\ref{cor:modg}
says that the model agrees with \(P,Q\) through every order below
\(g\).  Every partial model satisfies the reciprocal degree bounds:
indeed \(\deg_X[\tau^r]S^k\leq gk-r\), and the resulting sector
estimate is written explicitly at the end of this section.  Hence each
residual coefficient used below inherits the original degree bound.

Suppose inductively that the model agrees through order \(j-1\), where
\(j\geq g\), and write
\begin{equation}
\begin{aligned}
 P-P_{\mathrm{mod}}
 &=\tau^jT_j(X)+O(\tau^{j+1}),\\
 Q-Q_{\mathrm{mod}}
 &=\tau^jU_j(X)+O(\tau^{j+1}).
\end{aligned}
\label{eq:later-residual}
\end{equation}
At least one of \(T_j,U_j\) is nonzero.

\begin{lemma}[No strict root after the subresonant bridge]
\label{lem:no-strict-later}
If \(T_j\neq0\), put
\[
 \sigma_\alpha=\min\{\ord_\alpha T_j,a\},
 \qquad h_\alpha=a-\sigma_\alpha
\]
at each root of \(R\).  Then
\begin{equation}
 j\geq gh_\alpha
\label{eq:no-strict}
\end{equation}
for every \(\alpha\).  The analogous statement holds for \(U_j\), with
\(a\) replaced by \(b\).
\end{lemma}

\begin{proof}
Suppose \(j<gh_\alpha\) at some root.  In the coordinate
\(s=S(X,\tau)\), the point supplied by the residual lies below the
slope-\(g\) line from \(s^a\).  Hence the actual first local face is
strict.

Every positive term already in the \(P\)-model has the form
\[
 A_q\tau^{gq}s^{a-q}.
\]
For the strict primitive weights \(d',h'\), its excess over the weight
of \(s^a\) is
\[
 q(gh'-d')>0.
\]
Thus existing resonant sectors do not alter the first strict face.
Proposition~\ref{prop:strict-face} makes that face a common reduced
shift with primitive order \(1\leq\delta<g\).  It contains a nonzero
term \(a\kappa\tau^\delta s^{a-1}\).

But the model in the \(s\)-coordinate has positive orders only at
multiples of \(g\), while \eqref{eq:later-residual} has no residual
below \(j\).  Since \(\delta<g\leq j\), this is impossible.
The proof for a strict face detected from \(U_j\) is symmetric; the
common face would create the same forbidden primitive residual in
\(P\).
\end{proof}

\begin{proposition}[Classification of the least residual]
\label{prop:least-residual}
If \(T_j\neq0\), then for some integer \(q\) with \(1\leq q\leq a\)
and some \(c_q\in\C^\times\),
\begin{equation}
 j=gq,\qquad T_j=c_qR^{a-q}.
\label{eq:psector}
\end{equation}
If \(U_j\neq0\), then for some \(1\leq q\leq b\) and
\(d_q\in\C^\times\),
\begin{equation}
 j=gq,\qquad U_j=d_qR^{b-q}.
\label{eq:qsector}
\end{equation}
\end{proposition}

\begin{proof}
Assume \(T_j\neq0\).  Squarefreeness, the reciprocal bound, and the
truncated orders in Lemma~\ref{lem:no-strict-later} give
\[
 \sum_\alpha\sigma_\alpha
 \leq\deg T_j\leq ga-j,
\qquad\hbox{so}\qquad
 j\leq\sum_\alpha h_\alpha.
\]
On the other hand \eqref{eq:no-strict} gives
\[
 \sum_\alpha h_\alpha\leq j.
\]
All inequalities are equalities.  There are exactly \(g\) roots and
every \(h_\alpha\leq j/g\); equality of the sum forces
\[
 j=gq,\qquad h_\alpha=q\quad\hbox{for every }\alpha.
\]
The truncation is now inactive, \(1\leq q\leq a\), and
\[
 \ord_\alpha T_j=a-q
\]
at every root.  Its degree saturates the bound:
\[
 \deg T_j=g(a-q)=n-j.
\]
There is no room for an additional zero, so
\(T_j=c_qR^{a-q}\).  The proof for \(U_j\) is identical.
\end{proof}

\begin{proof}[Completion of the normal-form induction]
At order \(j=gq\),
\[
 [\tau^{gq}]
 \bigl(\tau^{gq}S^{a-q}\bigr)=R^{a-q}.
\]
Thus replacing \(A_q\) by \(A_q+c_q\) absorbs
\eqref{eq:psector} without changing an earlier coefficient.  Similarly
\(B_q\mapsto B_q+d_q\) absorbs \eqref{eq:qsector}.  If both residuals
are present, the changes are made simultaneously.  If only one is
present, its own scalar is changed; no relation between \(A\) and \(B\)
is imposed.

The induction terminates because \(P,Q\) are polynomials and their
reciprocal bounds give \(\tau\)-degrees at most \(n,m\).  It proves
\eqref{eq:normal-form}.

It remains to check that the model itself stays reciprocal.  From
\eqref{eq:S},
\[
 \deg_X[\tau^r]S^k\leq gk-r.
\]
Therefore
\[
 \deg_X[\tau^r]
 \bigl(\tau^{gq}S^{a-q}\bigr)
 \leq g(a-q)-(r-gq)=n-r,
\]
and analogously for \(Q\).  Also
\[
 gq+(g-1)(a-q)\leq n,
\]
so no excessive \(\tau\)-degree is introduced.  This completes the
proof of the asserted exact normal form in
Theorem~\ref{thm:normal-form}.
\end{proof}

\section{The resonant cone has zero bracket}
\label{sec:zero}

\begin{proposition}[Exact zero-bracket identity]
\label{prop:zero}
Let \(S\in\C[X,\tau]\) be arbitrary and put
\[
 z=\frac{\tau^g}{S},\qquad
 f(S,\tau)=S^aA(z),\qquad
 k(S,\tau)=S^bB(z),
\]
where \(A,B\) are arbitrary polynomials for which the displayed
expressions are polynomial.  If \(n=ga\) and \(m=gb\), then
\[
 \KS(f,k)=0,\qquad \KX(f(S,\tau),k(S,\tau))=0.
\]
\end{proposition}

\begin{proof}
At fixed \(S\),
\[
\begin{aligned}
 f_S&=S^{a-1}(aA-zA'),&
 f_\tau&=gS^azA'/\tau,\\
 k_S&=S^{b-1}(bB-zB'),&
 k_\tau&=gS^bzB'/\tau.
\end{aligned}
\]
Substitution into
\[
 \KS(f,k)
 =\tau(f_Sk_\tau-f_\tau k_S)+gafk_S-gbkf_S
\]
gives zero after cancellation, with no condition relating \(A\) and
\(B\).  Lemma~\ref{lem:covariance} then gives the \(X\)-coordinate
identity.
\end{proof}

\begin{proof}[Proof of Theorem~\ref{thm:main}]
The strict-face, subresonant-gluing, and least-residual results above
force the exact form \eqref{eq:normal-form}.
Proposition~\ref{prop:zero} then gives
\[
 \KX(P,Q)=0,
\]
contradicting
\[
 \KX(P,Q)=-c\tau^{n+m-2},\qquad c\neq0.
\]
\end{proof}

\begin{remark}[The case \(g=2\)]
There is no exceptional squarefree case at \(g=2\).
The gap \eqref{eq:weight-gap} remains positive, the subresonant bridge
handles the single order \(1\), the truncated degree average is
unchanged, and Proposition~\ref{prop:zero} is exact for every
\(g\geq1\).
\end{remark}

\section{Why repeated roots remain}
\label{sec:repeated}

Squarefreeness enters the proof in two indispensable ways.

\begin{enumerate}[label=(\roman*)]
\item At a simple root, \(s=S(X,\tau)\) is a formal local coordinate
  and \(S_X\) is a unit, so the bracket covariance
  \eqref{eq:covariance} preserves the initial face.
\item The \(g\) distinct roots turn reciprocal degree bounds into the
  sharp average used in Lemma~\ref{lem:strict-root} and
  Proposition~\ref{prop:least-residual}.
\end{enumerate}

If a root of \(R\) has multiplicity \(e\geq2\), choose instead a local
coordinate with \(R=s^e\).  Put
\[
 p=ae,\qquad q=be,\qquad \gamma=\frac ge.
\]
For a strict face
\[
 P_0=s^pA(z),\qquad z=\frac{\tau^{d'}}{s^{h'}},
\qquad \frac{d'}{h'}<\frac ge,
\]
the matched bracket is
\[
 (d'-\gamma h')s^{p+q-1}z
 \bigl(pAB'-qA'B\bigr).
\]
When this face lies below the scalar forcing, it follows only that
\[
 A=C^a,\qquad B=C^b,
\]
and the common local factor is
\begin{equation}
 s^eC\!\left(\frac{\tau^{d'}}{s^{h'}}\right).
\label{eq:repeated-factor}
\end{equation}
Polynomiality allows \(\deg C>1\), so
\eqref{eq:repeated-factor} can split the repeated factor.  A separate
degree argument rules out an apparent strict face whose weight exactly
equals the scalar forcing, but it does not make \(C\) linear.

There is also a global first-layer kernel.  Write
\[
 R=\prod_i(X-\alpha_i)^{e_i}
\]
and define the Kummer period
\[
 \kappa_R
 =\operatorname{lcm}_i\frac{g}{\gcd(g,e_i)}.
\]
At a fixed base \(R\), before the first nonzero \(P\)-coefficient (or
after subtracting the common contribution at that coefficient), the
homogeneous equation for a \(Q\)-coefficient at order \(j\) is
\[
 gR\,U_j'+(j-m)R'U_j=0.
\]
For an actually occupied coefficient, so \(0\leq j\leq m\), it has a
nonzero polynomial solution precisely when \(\kappa_R\mid j\), in which
case
\[
 U_j=cR^{\,b-j/g}.
\]
Thus repeated roots can support sub-\(g\) Kummer modes.  Even when the
globally first \(P\)-layer is a common factor deformation, iterating it
can change the factor multiplicities and the next Newton support.
The covariance identity alone does not supply the missing secondary
Newton--Rees induction.

Accordingly the following are not consequences of this paper:

\begin{itemize}
\item a normal form modulo \(\tau^g\) for an arbitrary repeated-root
  boundary;
\item elimination of the resonant slope \(g/e\);
\item elimination of Kummer modes after nonlinear interaction; or
\item a proof of the plane Jacobian conjecture.
\end{itemize}

\section{Reproducibility and checking}
\label{sec:repro}

The proof is symbolic, but its recurrent identities and integer
bookkeeping have exact companion checks in the repository containing
this source.  From the repository root, run the following scripts with
\texttt{python3} in an environment containing SymPy:
\begingroup\footnotesize
\begin{itemize}
\item \nolinkurl{current_context/verify_standard_system_single_binomial_face_no_scalar.py}
\item \nolinkurl{current_context/verify_standard_system_strict_compact_face_bridge.py}
\item \nolinkurl{current_context/verify_standard_system_squarefree_common_root_exclusion.py}
\end{itemize}
\endgroup
The first script checks the
telescoping coefficients and endpoint arithmetic in
Lemma~\ref{lem:binomial}.  The second checks
\eqref{eq:coset}, \eqref{eq:weight-gap}, coordinate covariance, and
the gluing degree arithmetic.  The third checks the strict-face formula,
the fact that every existing \(g\)-sector lies above a strict face, the
truncated degree-average lemma, sector absorption, reciprocal bounds,
and the exact zero-bracket identity.

For the scope statements in Section~\ref{sec:repeated}, the companion
checks are
\begingroup\footnotesize
\begin{itemize}
\item \nolinkurl{current_context/verify_standard_system_repeated_root_compact_face.py}
\item \nolinkurl{current_context/verify_standard_system_repeated_root_global_first_layer.py}
\item \nolinkurl{current_context/verify_standard_system_repeated_root_kummer_nonlinear_audit.py}
\end{itemize}
\endgroup
These programs use exact rational and polynomial arithmetic.  They are
not probabilistic and do not use floating point.  They are also not a
substitute for the Newton-face proof: in particular, automatic face
matching, local-to-global divisibility, and the logical use of
least-residual minimality are mathematical arguments in the text.

\paragraph{Trust boundary.}
The main theorem has no computer-exhaustive hypothesis.  A reader can
verify every displayed identity by differentiation and every discrete
step by the stated divisibility and degree inequalities.  The scripts
serve as regression checks against sign, weight, and indexing errors.

\paragraph{Disclosure.}
AI systems assisted with symbolic-search orchestration, adversarial
review, exact-code checking, and preparation of this draft.  The named
author is responsible for the mathematical claims, source, and
independent verification before dissemination.

\begin{thebibliography}{9}

\bibitem{BiggersMohFried}
R.~Biggers, T.~T. Moh, and M.~Fried,
\emph{On the Jacobian conjecture and the configurations of roots},
J. Reine Angew. Math. \textbf{340} (1983), 140--213,
\url{https://doi.org/10.1515/crll.1983.340.140}.

\bibitem{GGVSystem}
J.~A. Guccione, J.~J. Guccione, and C.~Valqui,
\emph{A system of polynomial equations related to the Jacobian
Conjecture},
preprint, arXiv:1406.0886, 2014,
\url{https://arxiv.org/abs/1406.0886}.

\bibitem{GGVLowerSide}
J.~A. Guccione, J.~J. Guccione, and C.~Valqui,
\emph{The two-dimensional Jacobian conjecture and the lower side of the
Newton polygon},
preprint, arXiv:1605.09430, 2016,
\url{https://arxiv.org/abs/1605.09430}.

\bibitem{GGHVAlgorithms}
J.~A. Guccione, J.~J. Guccione, R.~Horruitiner, and C.~Valqui,
\emph{Some algorithms related to the Jacobian Conjecture},
preprint, arXiv:1708.07936, 2017,
\url{https://arxiv.org/abs/1708.07936}.

\bibitem{GGHV2022}
J.~A. Guccione, J.~J. Guccione, R.~Horruitiner, and C.~Valqui,
\emph{Increasing the degree of a possible counterexample to the
Jacobian Conjecture from 100 to 108},
preprint, arXiv:2204.14178, 2022,
\url{https://arxiv.org/abs/2204.14178}.

\bibitem{MakarLimanovTrakhtenberg}
L.~Makar-Limanov and L.~Trakhtenberg,
\emph{Properties of a Jacobian mate},
S\~{a}o Paulo J. Math. Sci. \textbf{20} (2026), Article 16,
\url{https://doi.org/10.1007/s40863-025-00520-4};
open preprint MPIM 2024(33),
\url{https://archive.mpim-bonn.mpg.de/id/eprint/5148/}.

\bibitem{Shaska}
T.~Shaska,
\emph{Graded Keller maps and the Jacobian Conjecture},
preprint, arXiv:2607.20210v1, 2026,
\url{https://arxiv.org/abs/2607.20210}.

\end{thebibliography}

\end{document}
